

% 04 Apr: AI clean
The Congregation.

II.


\subsection{Installment 2: “Folkebladet” May 1, 1918}

A little further concerning our statements regarding the congregation.
We have said that a congregation exists wherever two or three are gathered in Jesus’ name, or, in other words, where two or more believers are gathered around the Word of God and the sacraments. This is Luther’s view, and this is, it seems to us, in accordance with article 2 in *Living Principles*, which says: “The congregation consists of believing people, who by the use of the means of grace and the spiritual gifts, according to the direction of God’s Word, seek their own salvation and eternal blessedness and that of all their fellow men.”

We believe, in accordance with point 3, that according to the New Testament the congregation requires an outward organization, with a register of its members, the election of officers, fixed times and places for its gatherings, etc.

The congregation is, however, Christ’s body, and one becomes a member of His body only by a true and living faith in Jesus Christ.

By nature we are all dead; but by a new birth, wrought by the Holy Spirit, we receive life in God, become one with Him, and thus true members of the congregation.

We say that the congregation is the true form of the Kingdom of God on earth; but we cannot simply say that when someone becomes a member of the congregation’s outward organization, he thereby also becomes a member of the Kingdom of God; for membership in the Kingdom of God depends upon living union with God. And one can have living fellowship with God without being a member of the congregation’s outward organization. In accordance with this understanding, we also use the word “congregation” to denote the gathering of all believers wherever they are found, within or outside an organization. For in the third article we confess that the church (the congregation) is one; but this unity consists not in organization, not in doctrinal confession, not in liturgical form, but in life, life in God.

Here, then, a distinction is made between the congregation according to its inward essence and according to its outward form. Such a distinction is in full accordance both with the Word of God and with *Living Principles*. When, however, the congregation according to its inward, spiritual essence has been termed “the invisible church,” this is not a happy expression, since the Word of God plainly says that it is to be as the city set on the hill that cannot be hidden. To be sure, the Savior says: “The Kingdom of God does not come in such a way as to be observed; neither shall they say, See here, or, See there! for behold, the Kingdom of God is within you.” But He also says that believers are to be light and salt in the world, that the light is not to be put under the bushel, but on the stand, so that it may shine for all in the house.

“Let your light so shine before people, that they may see your good works and praise your Father in heaven.”

These words of the Savior may be applied to the congregation as a whole just as well as to the individual believers; for the congregation becomes such as its individual members are. If the individuals are shining lights, then the congregation also will shine. If the life is warm and vigorous in the individuals, it will also manifest itself in the congregation and outward from it. Our life is hidden with Christ in God; but it must also come visibly forth in the congregation, in the individual members. Believers are to be living epistles, known and read by all. God’s Spirit is indeed invisible, but He makes Himself known by His working in and through believers. When they walk worthy of their calling, it will always be possible to say of them: they have been with Jesus; they have been in the school of the Holy Spirit. For them God’s will is always law; they walk with the Lord, they serve Him.

Life itself is invisible; but it manifests itself by the forms it assumes, by the development through which it passes. Where there is life, there is movement, activity, labor. Stagnation is death.

The disciples of Jesus Christ are called to labor. They have received life not only to enjoy, but also to give. They are to bear fruit for the Kingdom of God; indeed, they are even to receive grace to suffer with Christ, to fill up in their flesh what is yet lacking in Christ’s afflictions, for His body, which is the congregation.

How much affliction are we willing to bear for God’s congregation? How great sacrifices are we willing to make? How much labor are we willing to perform? Is it a matter of the heart for us to do our part that the Kingdom of God may flourish among us? Or are our fair words about the congregation only empty phrases? May God give us more and more of the mind of Jesus Christ, that we may become more and more like Him in serving, suffering, sacrificing.

Self-sacrificing labor is required. The individual believer is to labor for the advancement of the Kingdom of God; the whole congregation is to be a well-organized working body. The gifts of grace must be cultivated and put to use. It is not God’s thought that believers should sit still and inactive, each by himself, one here and one there; He wills that they should come together in order to encourage, urge, and help one another. They need such mutual help both as regards their own salvation and the salvation of other souls. Besides, brotherly love will draw God’s children together; spiritual kinship, the need for common edification and prayer, will unite them, make them one. The believers in the first congregation remained steadfast in the apostles’ teaching and in the fellowship, in the breaking of bread and in the prayers. And these are at all times the true marks of a sound congregational life. There is therefore something amiss with those who do not feel the need to gather with others around the Word of God and the sacraments for common edification, common prayer, and common labor. They lack, at any rate, the proper congregational sense, the proper sense of fellowship, and the proper responsibility toward the advancement of the Kingdom of God on earth.

A congregational organization is also needed so that there may be order and discipline and so that with united strength one may carry out work that the individual cannot accomplish alone.

But it ought always to be remembered that the congregation properly consists of believing people, that it is to be holy and blameless, without spot or wrinkle, and it ought to be our aim, with God’s help, to present the congregation in its proper form according to the plan the Lord has indicated in His Word. Here is our task. It is a high and glorious goal. We shall scarcely attain it fully here upon earth; but to come as near to it as possible ought to be our striving.

To labor unceasingly for the congregation’s quickening and liberation among us is our task. Awakening, Spirit, and life are the essential thing. It is not enough to presuppose life; we must do our part that there may be life, more and more life, among us.

Yet awakening to life does not come merely by speaking about awakening and crying out for awakening.

Our stock expressions about awakening, Spirit, and life can, like so much else, easily become customary forms of speech and empty phrases without the slightest effect.

To stir up a religious movement with much fuss and display, moving stories, many tears, and appeal to the emotions is not, without more, true awakening either.

All this may be very harmful.

God’s Spirit alone can awaken to life.

Human contrivances are of no use.

We must therefore let God’s Spirit do His work in us and through us.

We must become fervent in spirit, steadfast in prayer, and zealous for every good work.

It must be our highest delight and longing to be the Spirit’s willing instruments in bearing forth the testimony of salvation and grace in Jesus Christ and leading sinners to Him.

But it is true that souls can be saved one by one and liberated from the power and dominion of sin without the congregation to which they belong being liberated. Therefore we say with Ditzel that “next to the question of the salvation of the individual soul there is no more important question than that of the congregation.” And therefore we continue to say with point 2 in the Free Church’s rules that “its (the Lutheran Free Church’s) purpose is to labor for the quickening and liberation of all Norwegian Lutheran congregations, so that according to their calling and their ability they may work in spiritual independence and freedom for the cause of the Kingdom of God at home and abroad through such agencies and institutions as the congregations themselves may determine.”

So then, it is both quickening and liberation. Of spiritual liberation there can be no question apart from spiritual life. Freedom must come as a fruit of the Spirit; for “where the Spirit of the Lord is, there is freedom.”
